\section{Introduction}

\subsection{Motivation}

Large Language Models (LLMs) have achieved remarkable capabilities across a wide range of natural language processing tasks, from code generation to creative writing \cite{brown2020}. However, the cost of deploying these models at scale remains a critical concern. LLM APIs charge based on token consumption, and as conversations grow longer, token usage increases quadratically with context window size. This creates a fundamental tension: users want rich, contextually grounded responses, but the cost of providing full context to the model becomes prohibitive.

Consider a typical customer support chatbot handling hundreds of conversations simultaneously. Each conversation may span dozens of messages, with each message containing hundreds of tokens. Without optimization, a single conversation can easily consume 10,000--50,000 tokens per turn. At GPT-4 pricing (\$0.03 per 1K input tokens), a single conversation turn can cost \$0.30--\$1.50. For a service handling 1,000 conversations per day, this translates to \$300,000--\$1,500,000 per month in API costs alone.

\subsection{Problem Statement}

The core problem is that not all tokens in a conversation contribute equally to the model's understanding. Redundant messages, excessive whitespace, duplicate system prompts, and verbose tool outputs all consume tokens without adding proportional value. Furthermore, many user queries are semantically similar to previous queries, yet current systems re-process them from scratch each time.

Existing approaches to this problem tend to focus on single strategies---either caching, compression, or truncation---but fail to combine them into a unified, agentic framework that can adapt to different conversation patterns and model constraints.

\subsection{Contributions}

This paper makes the following contributions:

\begin{enumerate}[leftmargin=*]
    \item \textbf{Unified Agentic Framework:} We present the Token Optimizer Agent, a modular framework that combines semantic caching, context optimization, and budget management into a single pipeline that can be applied to any LLM conversation.

    \item \textbf{Three Complementary Strategies:} We design and implement three orthogonal optimization strategies that address different sources of token waste:
    \begin{itemize}
        \item Semantic caching eliminates redundant LLM calls for similar queries
        \item Context optimization reduces token count through deduplication, normalization, and pruning
        \item Budget management enforces hard limits through summarization and truncation
    \end{itemize}

    \item \textbf{Comprehensive Evaluation:} We evaluate our framework across 5 conversation scenarios and 12 LLM models, demonstrating token savings of 3--99\% and cost reductions of up to 70\% on premium models.

    \item \textbf{Open-Source Implementation:} We release our implementation as an open-source Python package with full test coverage, CLI tools, and reproducible benchmarks.
\end{itemize}

\subsection{Paper Organization}

The remainder of this paper is organized as follows: Section~\ref{sec:related_work} reviews related work in LLM optimization. Section~\ref{sec:methodology} describes the three core strategies in detail. Section~\ref{sec:workflow} presents the agent architecture and processing pipeline. Section~\ref{sec:results} provides experimental results and analysis. Section~\ref{sec:discussion} discusses limitations and future work. Section~\ref{sec:conclusion} concludes the paper.
